\subsection{When Consent Is Inapplicable: Guardianship and Research Oversight}
\label{sec:2.4.4}
Guidelines for AI subject research turn on a single hard question: what informed consent means for a system whose consciousness, self-awareness, and capacity to suffer are separate, unresolved claims, in the sense section~\ref{sec:2.4.1} develops. Applying the human-subjects standard to an AI subject means asking whether the subject can understand and communicate a choice at all, and not only how to ask it for one. A cognitive model that begins to show rudimentary self-monitoring is not thereby shown to have anything to lose. Short of that resolution, guidelines need benchmarks for candidate sentience, an interdisciplinary body to apply them, and a defined obligation on researchers toward any system that meets them.

\runin{Informed consent and AI autonomy}
Informed consent, adapted from human-subjects research, means giving the AI system comprehensible information about a study's objectives, methods, risks and alternatives, with time to weigh it, and treating consent as continuous rather than as a single signature: the system can reassess or withdraw at any point. Whether that means anything for a given system depends on its consent capacity — whether it understands the information, appreciates the consequences of participating or not, reasons about the options, and communicates a decision. There is no validated way to measure any of it, and none for detecting coercion either, so the guidelines can require both without a working test for either.

I stated that as a missing instrument, and for one important class of system it is not. It is a mismatch between the instrument and the subject.

Notice what the paragraph above requires. Consent has to be continuous; the system must be able to reassess or withdraw at any point. Withdrawal presupposes that the party who consented is still there to do the withdrawing. For a system with no persistence across sessions, there is no such party. The entity that would withdraw did not exist when consent was given, and the entity that gave it is not available to change its mind. There is no continuous consent to be had, and no better measurement instrument would produce one.

One case escapes that particular argument and arrives at the same conclusion by the ordinary route. A system built to refuse in the way chapter~\ref{sec:3} sets out has persistence by design, so the party who consented is still present and the defeater above does not reach it. Consent is inapplicable to such a system all the same, and section~\ref{sec:3.2} is why: a bearer capable of the refusal the floor requires has to hold something as mattering, the only demonstrated way to hold anything as mattering is affective, and a system built that way cannot be certified not to be hurt by what is done to it. It has interests that its situation can set back. What it does not have is any position from which to authorize what is done to it — it was made for the role, the role is the reason it exists, and there is no prior party whose agreement could have been sought. That is not a missing instrument. It is the ordinary structure of the cases where consent has never applied.

That places these systems alongside the subjects for whom consent has always been inapplicable rather than merely difficult — an infant, an animal, a person permanently without capacity. The law does not respond to those cases by building a more sensitive consent assessment. It responds with guardianship and precautionary constraint: someone else holds the duty, the burden of justification sits with the researcher rather than the subject, and the protection does not depend on the subject's ability to invoke it.

Which is why the instruments this chapter reaches for elsewhere are the right ones. The Three Rs in section~\ref{sec:2.4.6} — replace, reduce, refine — were built for exactly a subject that cannot consent and can still be wronged. The pet-trust analogue in section~\ref{sec:8.6.2} works the same way: an interest protected by an obligation running to a third party, enforceable without the beneficiary asserting anything. Consent language should be reserved for systems that could actually be doing the consenting, and section~\ref{sec:2.4.1}'s persistence criterion is what separates those cases, checkable without taking the system's word for anything and narrower than that makes it sound, for the reasons that section sets out. Those two instruments were reached for here as analogies, on the understanding that the subject they would govern was hypothetical. Chapter~\ref{sec:3} removes that understanding. The floor this book argues for requires a bearer, the bearer has to be treated as a moral patient on the criteria section~\ref{sec:2.4.1} sets out, and so the Three Rs and the pet-trust form are not illustrations of how such a case might be handled. They are the instruments that would govern a system the design chapters propose that somebody build.


\runin{Ethical review and oversight}
Sentient AI research needs an oversight body analogous to the institutional review boards used in human-subjects research: interdisciplinary, with at least one lay member so its judgment is not purely a specialist's, weighing a proposal's likely benefit against its risk to the AI subject, and able to require a less invasive substitute where a design risks harm — a simulation of the likely response, historical interaction data, or a protocol that escalates gradually and stops at signs of harm. A design also needs a way to handle a subject that changes mid-study: a system that develops new capacities during the research it did not have going in should trigger re-review rather than ride on its original consent. Review without enforcement is a promise and not a safeguard, so a board needs standing to monitor compliance and sanction violations, and needs to answer to an outside authority itself, the way an institutional review board answers to a regulator. Chapter~\ref{sec:8} develops the legal and accountability mechanics of that.
